\section{Исходная программа}

\subsection{Заголовочный файл monte\_carlo.h}
\begin{center}
\lstinputlisting[language=C,caption=Заголовочный файл с объявлениями,captionpos=b]{inc/monte_carlo.h}
\end{center}

\subsection{Основная программа monte\_carlo.c}
\begin{center}
\lstinputlisting[language=C,caption=Реализация многопоточного алгоритма Монте-Карло,captionpos=b]{inc/monte_carlo.c}
\end{center}


\subsection{Пример использования программы}
\begin{verbatim}
$ gcc -o monte_carlo monte_carlo.c -lm -lpthread
$ ./monte_carlo -h
Использование: ./monte_carlo -r RADIUS -p POINTS -t THREADS [опции]
Опции:
  -r RADIUS    Радиус окружности (по умолчанию: 1.0)
  -p POINTS    Общее количество точек (по умолчанию: 1000000)
  -t THREADS   Максимальное количество потоков (обязательно)
  -v           Подробный вывод
  -h           Эта справка
Пример: ./monte_carlo -r 5.0 -p 10000000 -t 8

$ ./monte_carlo -r 5.0 -p 10000000 -t 4
=== Вычисление площади окружности методом Монте-Карло ===
Радиус: 5.00
Количество точек: 10000000
Максимальное количество потоков: 4
Количество ядер процессора: 8

=== Результаты ===
Точек внутри окружности: 7853542 из 10000000 (78.54%)
Вычисленная площадь: 78.535420
Теоретическая площадь (π * R²): 78.539816
Погрешность: 0.01%
Время выполнения: 0.342 секунд
Скорость обработки: 29239766 точек/сек
Оценка числа π: 3.141417 (ошибка: 0.01%)
\end{verbatim}

\subsection{Демонстрация потоков в системе}
Для демонстрации количества потоков, используемых программой:
\begin{verbatim}
# Запуск программы в фоновом режиме
$ ./monte_carlo -r 5.0 -p 100000000 -t 8 &
[1] 12345

# Просмотр потоков процесса
$ ps -eLf | grep monte_carlo
user    12345 12344 12345  0 12:34 pts/0 00:00:00 ./monte_carlo -r 5.0 -p ...
user    12345 12344 12346  0 12:34 pts/0 00:00:00 ./monte_carlo -r 5.0 -p ...
user    12345 12344 12347  0 12:34 pts/0 00:00:00 ./monte_carlo -r 5.0 -p ...
user    12345 12344 12348  0 12:34 pts/0 00:00:00 ./monte_carlo -r 5.0 -p ...
... (всего 9 потоков: 8 рабочих + 1 основной)

# Использование top для мониторинга потоков
$ top -H -p 12345
\end{verbatim}

\subsection{Анализ системных вызовов с помощью strace}
\begin{verbatim}
$ strace -f -o trace.log ./monte_carlo -t 4 -p 1000000
$ grep -E "pthread_create|clone" trace.log
[pid 12345] clone(child_stack=NULL, flags=CLONE_THREAD|CLONE_SYSVSEM|CLONE_SETTLS
      |CLONE_PARENT_SETTID|CLONE_CHILD_CLEARTID, parent_tidptr=0x7f8a1d7d4a10,
      tls=0x7f8a1d7d4700, child_tidptr=0x7f8a1d7d4a10) = 12346
[pid 12345] clone(...) = 12347
[pid 12345] clone(...) = 12348
[pid 12345] clone(...) = 12349
\end{verbatim}